\documentclass[11pt,a4paper,spanish]{book}
\usepackage{unir}
\usepackage{apacite}
\setlength{\headheight}{26pt}
\usepackage{hyperref}
\numberwithin{equation}{chapter}
%\usepackage{caption}
\numberwithin{figure}{chapter}
%\usepackage{chngcntr}
%\counterwithin{equation}{chapter}
%\counterwithin{figure}{chapter}
%\counterwithin{table}{chapter}
%\renewcommand\theequation{\thechapter.\arabic{equation}}
%\renewcommand\thefigure{\thechapter.\arabic{figure}}
%\renewcommand\thetable{\thechapter.\arabic{table}}
%\renewcommand\theequation{\thechapter.\arabic{equation}}
%\counterwithin{figure}{chapter}
%\counterwithin{table}{chapter}
%\renewcommand\thefigure{\thechapter.\arabic{figure}}
%\renewcommand\thetable{\thechapter.\arabic{table}}
%\renewcommand\thefigure{\thechapter.\arabic{figure}}
%\makeatletter
%\renewcommand\p@figure{\thechapter..\arabic{figure}}
%\makeatother



%---------------------------
%título del trabajo y autor
%---------------------------
\title{Auge del Indie: Análisis Predictivo y Visualización del Mercado 
de Videojuegos Independientes}
\titulacion{Máster Universitario en Análisis y Visualización de Datos 
Masivos / Visual Analytics \& Big Data}
\author{David Vicente Moya}
\date{\today}
\director{Diego Fernández García}
\nombreciudad{Teruel}

%---------------------------
%marges
%---------------------------
%\usepackage[margin=1.9cm]{geometry}
%---------------------------
%---------------------------
%---------------------------
%---------------------------
\begin{document}
\renewcommand{\listfigurename}{Índice de Ilustraciones}
\renewcommand{\listtablename}{Índice de Tablas}
\renewcommand{\contentsname}{Índice de Contenidos}
\renewcommand{\figurename}{Figura}
\renewcommand{\tablename}{Tabla} 

\maketitle

\frontmatter
\tableofcontents
\listoffigures
\listoftables

\chapter{Resumen}
{\bf Nota:} En este apartado se introducirá un breve resumen en 
español del trabajo realizado (extensión máxima: 150 palabras). 
Este resumen debe incluir el objetivo o propósito de la investigación, 
la metodología, los resultados y las conclusiones.


{\bf Palabras Clave:} Se deben incluir de 3 a 5 palabras claves en 
español

\chapter{Abstract}
{\bf Nota:} En este apartado se introducirá un breve resumen en 
español del trabajo realizado (extensión máxima: 150 palabras). Este 
resumen debe incluir el objetivo o propósito de la investigación, la 
metodología, los resultados y las conclusiones.


{\bf Palabras Clave:} Se deben incluir de 3 a 5 palabras claves en 
inglés




\mainmatter
\chapter{Introducción}
\section{Motivación y planteamiento del problema}
\section{Pregunta de investigación}
\section{Estructura del documento}



\chapter{Contexto y Estado del Arte}

\section{Indie y AAA}
En la última década el sector del videojuego ha presenciado un 
crecimiento inesperado de los estudios independientes. La democratización
de herramientas de desarrollo (Unity, Godot, Unreal Engine), la distribución de 
copias digital a través de plataformas online como Steam o Itch.io, y el 
aumento de la popularidad de plataformas de micromecenazgo, han reducido las 
barreras de entrada en el mercado para estudios pequeños y desarrolladores 
solitarios. Títulos como \textit{Star Dew Valley}, \textit{Hollow Knight} o 
\textit{Hades} demuestran que un equipo reducido puede alcanzar un éxito comercial 
y comparable al de grandes producciones, llegando incluso a ocupar gran 
parte del mercado~\cite{informe_mercado_2024}.

\vspace{1em}
Este auge a quebrado el balance que existía previamente entre el desarrollo 
independiente y el denominado AAA. Mientras que este úlrimo se caracteriza por 
grandes presupuestos, equipos extensos y estrategias de marketing masivas, el 
desarrollo indie suele asumir menos recursos a cambio de mayor libertad creativa. 
Sin embargo, la frontera que separaba ambos se desvanece cada día que pasa, con 
estudios "AA" y editores que respaldan a pequeños desarrolladores. A día 
de hoy, jugadores de todo ámbito tienen opiniones varias sobre lo que se puede o 
no considerar un videojuego independiente~\cite{definicion_indie_jugon}.

\vspace{1em}
El crecimiento de la oferta ha traído consigo un problema de saturación. El 
número de lanzamientos anuales en Steam ha pasado de unos pocos cientos a varios 
miles por año. En 2026 hemos llegado a presenciar \textit{"720 nuevos lanzamientos 
en la plataforma de Valve sólo durante la semana del 24 al 30 de agosto de 
2026"}~\cite{aumento_lanzamientos_indie}. Esta saturación dificulta la 
visibilidad de cada título. Además, los ingresos se concentran en un porcentaje 
reducido de juegos, mientras que la mayoría obtiene resultados 
modestos~\cite{inf_mercado_2026} un desarrollador, lanzar un juego equivale a 
competir en un entorno donde ser descubierto es tan difícil como desarrollarlo.

\vspace{1em}
A esto se le suma la dificultad de medir el éxito comercial. Cada fuente de datos 
(Steam API, SteamSpy, IGDB, Kaggle) presenta limitaciones de cobertura y precisión 
que condicionan cualquier análisis. Pero podemos permitirno limitar los datos de 
nuestro análisis a los datos históricos de Steam, ya que ocupa casi la totalidad 
del mercado digital. Según la revista Vandal, hasta un 75\% de los ingresos 
de los estudios proviene exclusivamente de la plataforma de 
Valve~\cite{steam_monopolio}.

\vspace{1em}
Existen múltiples factores que podrían influir en el éxito de un título: género y 
etiquetas, precio, fecha de lanzamiento, idiomas, plataformas, comunidad previa o 
reseñas tempranas. No obstante, la relación entre ellos no está clara y suele 
tratarse de forma intuitiva más que basada en datos. Aunque existen estudios 
previos sobre predicción de éxito en videojuegos, como la utilización de Random 
Forest para predecir el éxito comercial~\cite{vivero2025}, son escasos los enfoques 
centrados en el sector indie que combinen un modelo predictivo con una herramienta 
de visualización interactiva accesible.

\vspace{1em}
En este contexto, surge la necesidad de analizar de forma sistemática los datos 
históricos de Steam para identificar qué factores se asocian al éxito comercial 
indie y cómo ha evolucionado su cuota de mercado frente al AAA. Este trabajo busca 
aportar tanto un modelo predictivo como un dashboard interactivo útiles para 
desarrolladores, editores e inversores.

\section{Estado del Arte}

\chapter{Identificación de Requisitos}

\chapter{Objetivos}

\chapter{Desarrollo del trabajo}

\chapter{Conclusiones y Trabajo Futuro}

%En la Ecuación \eqref{eq:eq1secCTF}
%
%
%\begin{equation}\label{eq:eq1secCTF}
%M=\begin{pmatrix}
%	m_{11}&m_{12}\\
%	m_{21}&m_{22}
%\end{pmatrix}
%\end{equation}
%
%En la siguiente Tabla \ref{tab:tab1secCTF}
%
%\begin{table}[h]
%\centering
%\begin{tabular}{|c|c|}
%	\hline
%	1 & 2 \\
%	\hline
%	22 & 11 \\
%	\hline
%\end{tabular}
%\caption{Tabla 1}
%\label{tab:tab1secCTF}
%\end{table}
%
%En la siguiente Figura \ref{fig:fig1secCTF}
%
%\begin{figure}[h]
%\includegraphics[width= 0.8\textwidth]{logo_unir}
%\caption{Logo Unir}
%\label{fig:fig1secCTF}
%\end{figure}
%
%\cite{PIMENTEL2016744} \cite{da_S_Bessa_2023}

%\begin{thebibliography}{a}
%\bibitem{etiqueta} \textsc{Autores},
%\textit{nombre referencia.}
%Información addicional
%\end{thebibliography}
\bibliographystyle{apacite}
\bibliography{bibliografia}

\appendix
\chapter{Apendices}

\end{document}





















